\begin{lemma}[Ray closure]
Let $F$ be a nontrivial front on $X$.  For every $n\in X$, the ray
$F_n$ is a front on $X/n$, and
\[
F=\bigl\{\{n\}\cup s\mid n\in X,\ s\in F_n\bigr\}.
\]
\end{lemma}
\begin{proof}
Fix $n\in X$.  We first record two consequences of the hypotheses that will
be used repeatedly.  Since $F\ne\{\emptyset\}$ and $F$ is prefix-free in
the sense of \noderef{Front}, every member of $F$ is nonempty.  Indeed, if
$\emptyset\in F$, then $\mathsf{InitialSegment}(\emptyset,v)$ holds for
every finite $v$: it is equality when $v=\emptyset$, and otherwise take the
least element of $v$ in the second clause of \noderef{InitialSegment}.
Prefix-freeness would then give $v=\emptyset$ for every $v\in F$.  Density
from \noderef{Front}, applied to the infinite set $X$, makes $F$ nonempty,
so this would imply $F=\{\emptyset\}$, contrary to the hypothesis.  The
base disjunction in \noderef{Front} therefore yields
\[
\mathsf{FrontBase}(F)=X.
\tag{1}
\]

We prove that the ray is dense on $\mathsf{TailSet}(X,n)$.  Let $Y$ be an
infinite subset of $\mathsf{TailSet}(X,n)$, with the meaning of this tail
given by \noderef{TailSet}.  Every element of $Y$ belongs to $X$ and is
greater than $n$.  Consequently $n\notin Y$, the set $\{n\}\cup Y$ is
infinite (if it were finite, its subset $Y$ would be finite), and
$\{n\}\cup Y\subseteq X$: the new element $n$ is in $X$, and all other
elements are in $Y\subseteq X$.  By the density clause of \noderef{Front},
there are $u\in F$ and a witness $m\in\{n\}\cup Y$ such that
\[
 u=\{k\in\{n\}\cup Y\mid k<m\};
\]
this is exactly $\mathsf{ProperPrefixSet}(u,\{n\}\cup Y)$ by
\noderef{ProperPrefixSet}.  The preliminary observation gives
$u\ne\emptyset$.  Since $n$ is the least element of $\{n\}\cup Y$, the
witness $m$ is not $n$, so $m\in Y$ and $n\in u$.  Put
$s=u\setminus\{n\}$.  Then every element of $s$ is greater than $n$,
$\{n\}\cup s=u\in F$, and
\[
 s=\{k\in Y\mid k<m\},
\]
so $\mathsf{ProperPrefixSet}(s,Y)$ by \noderef{ProperPrefixSet}.  The
definition of the ray in \noderef{Ray} now gives $s\in F_n$.  Thus
$F_n$ is dense on $\mathsf{TailSet}(X,n)$.

The tail is infinite because $X$ is infinite by \noderef{Front}, while only
the finitely many naturals at most $n$ have been removed; hence
$\mathsf{TailSet}(X,n)$ satisfies the infinite-base clause of \noderef{Front}.
It remains to verify the base and prefix clauses.  If
$F_n=\{\emptyset\}$, its base clause holds by the first alternative in
\noderef{Front}.  Suppose instead that $F_n\ne\{\emptyset\}$.  We prove
both inclusions
\[
\mathsf{FrontBase}(F_n)=\mathsf{TailSet}(X,n).
\tag{2}
\]
For the inclusion from left to right, let $k\in\mathsf{FrontBase}(F_n)$.
By \noderef{FrontBase}, choose $s\in F_n$ with $k\in s$.  The ray
definition in \noderef{Ray} gives $n<k$ and
$\{n\}\cup s\in F$.  Therefore $k\in\mathsf{FrontBase}(F)=X$ by (1),
and hence $k\in\mathsf{TailSet}(X,n)$ by \noderef{TailSet}.

Conversely, let $k\in\mathsf{TailSet}(X,n)$.  Let
\[
 Y_k=\{k\}\cup\mathsf{TailSet}(X,k).
\]
The same finite-removal argument shows that $\mathsf{TailSet}(X,k)$ is
infinite, so $Y_k$ is infinite.  Also $Y_k\subseteq\mathsf{TailSet}(X,n)$:
the element $k$ lies in $X$ and satisfies $n<k$, while every element of the
second summand lies in $X$ and is greater than $k$.  Ray density, already
proved above, supplies $s\in F_n$ with
$\mathsf{ProperPrefixSet}(s,Y_k)$, using \noderef{ProperPrefixSet}.  This
$s$ cannot be empty.  For if $\emptyset\in F_n$, then
$\mathsf{InitialSegment}(\emptyset,t)$ holds for every $t$, and the
prefix-free argument using \noderef{InitialSegment} and the prefix clause
of \noderef{Front} would force every member of $F_n$ to be empty, contrary
to $F_n\ne\{\emptyset\}$.  A nonempty proper prefix of $Y_k$, whose least
element is $k$, contains $k$.  Thus $k\in s$, and \noderef{FrontBase} gives
$k\in\mathsf{FrontBase}(F_n)$.  This proves (2), so the base clause holds in
the nontrivial-ray case as well.

For prefix-freeness, take $s,t\in F_n$ and assume
$\mathsf{InitialSegment}(s,t)$.  Every element of $s$ and $t$ is greater
than $n$ by \noderef{Ray}, so $n$ is the common least element of
$\{n\}\cup s$ and $\{n\}\cup t$.  Inserting this common least element
preserves and reflects initial segments: directly from the two alternatives
in the definition of \noderef{InitialSegment},
\[
\mathsf{InitialSegment}(s,t)\quad\Longleftrightarrow\quad
\mathsf{InitialSegment}(\{n\}\cup s,\{n\}\cup t).
\]
The ray definition puts both inserted sets in $F$, so prefix-freeness of
\noderef{Front} gives $\{n\}\cup s=\{n\}\cup t$, and cancellation of the
common element gives $s=t$.  We have now verified all clauses of
\noderef{Front}, and therefore $F_n$ is a front on
$\mathsf{TailSet}(X,n)$.

It remains to prove the displayed decomposition.  For the inclusion from
left to right, let $u\in F$.  The preliminary observation says that $u$ is
nonempty; let $n$ be its least element.  Since
$n\in\mathsf{FrontBase}(F)=X$ by (1), put $s=u\setminus\{n\}$.  Every
element of $s$ is greater than $n$, and $\{n\}\cup s=u\in F$, so
$s\in F_n$ by \noderef{Ray}.  Hence
\[
u=\{n\}\cup s\in
\{v\mid\exists n\in X\;\exists s\in F_n\;v=\{n\}\cup s\}.
\]
For the reverse inclusion, let $n\in X$ and $s\in F_n$.  The definition of
\noderef{Ray} immediately says that $\{n\}\cup s\in F$.  Thus both
inclusions hold, proving the displayed equality and the lemma.
\end{proof}
